\section*{INTRODUCCIÓN}
\addcontentsline{toc}{section}{INTRODUCCIÓN}

La operación de una tienda minorista de electrónica requiere mantener coordinados el catálogo de productos, las existencias, las recepciones de mercancía, las ventas, la información de compradores y el trabajo de sus operadores. Cuando estos datos se actualizan por separado o de forma manual, resulta difícil conservar un estado común, localizar productos y mantener un registro de los movimientos. Tech-Flow aborda ese problema mediante una aplicación de escritorio que reúne esas tareas en una interfaz con acceso por roles y persistencia local en archivos CSV.

El sistema implementado utiliza Java 17 y JavaFX. El administrador gestiona la caja y supervisa o mantiene registros; el cajero prepara cotizaciones y ventas; el almacenero consulta el inventario y registra recepciones. Las operaciones comparten archivos locales, por lo que un cambio en el stock puede reflejarse entre ventas y entradas. No se encontró una base de datos relacional, un servicio web ni sincronización entre equipos: el alcance es una aplicación de escritorio que opera con los archivos disponibles en su directorio de datos.

El análisis técnico resulta relevante porque el proyecto combina varias preocupaciones: estructuras de datos en memoria para filas, resultados y carritos; archivos CSV como persistencia; clases de dominio que ejecutan reglas de negocio; pruebas automatizadas; y vistas JavaFX enlazadas con controladores. Revisar cada componente por separado y sus puntos de integración permite distinguir las responsabilidades que sí están separadas de aquellas que todavía se encuentran acopladas.

El contexto académico se evidencia en los datos de portada compartidos por el propio proyecto. Sin embargo, el repositorio no conserva una especificación original completa ni evidencia suficiente para reconstruir cada decisión de su concepción inicial. Por esa razón, este informe describe la evolución que puede verificarse en el código actual, en el historial Git visible y en documentos históricos consultados durante el análisis; cuando una afirmación procede únicamente de un registro anterior, se identifica como tal y no como una validación independiente.

La exposición se organiza en tres fases de análisis: definición y diseño de la estructura inicial; pruebas, refactorización y calidad; y arquitectura MVC, GUI y escenarios de validación. Esta organización sirve para explicar de forma lógica los resultados, pero no constituye prueba de que el equipo aplicara formalmente un proceso secuencial o una metodología concreta. El informe prioriza el estado actual del código y utiliza las capturas como evidencia visual de estados particulares, no como sustituto de pruebas funcionales.
